Several directions remain for the next phase of the project:

\textbf{Simulation Integration.}
The debate system currently operates on static sample observations. Our next step is to integrate it with the team's market simulation environment so that debate-generated trades are executed against realistic price dynamics, enabling us to measure PnL, Sharpe ratio, and maximum drawdown alongside reasoning quality.

\textbf{Quantitative Reasoning Evaluation.}
We plan to score the causal claims produced during debates using a T\textsuperscript{3}/CRIT-style rubric grounded in Pearl's causal hierarchy. Each claim is already tagged with a Pearl level (L1/L2/L3); the next step is to build an automated evaluator that checks whether the stated variables, assumptions, and falsifiers are internally consistent and whether L2/L3 claims actually demonstrate interventional or counterfactual reasoning rather than dressed-up associations.

\textbf{Correlation Analysis (RQ2).}
Once simulation integration is complete, we will perform joint analysis between reasoning quality scores and trading performance metrics to test whether improved causal reasoning correlates with better risk-adjusted returns.

\textbf{Failure Mode Analysis (RQ3).}
We will systematically compare outputs across agreeableness settings (Configs 4 vs.\ 5) and adversarial vs.\ non-adversarial setups (Configs 1 vs.\ 3) to measure whether structured debate reduces overconfidence, sycophancy, hallucinated causality, and groupthink. The agreeableness knob provides a controlled axis for studying sycophancy in particular.

\textbf{Distillation to Single-Agent Prompts.}
If debate produces measurably better reasoning, we plan to distill the most effective constraints (e.g., mandatory falsifiers, L2/L3 requirements, self-critique) back into single-agent prompts to test whether the benefits of multi-agent deliberation can be captured without the computational overhead of running multiple agents.

\textbf{Scaling and Robustness.}
We intend to test the system across a wider universe of stocks, longer time horizons, and different LLM backends to assess whether the observed patterns generalize beyond our initial experimental setup.
